%%%%%%%%%%%%%%%%%%%%%%%%%%%%%%%%%%%%%%%%%%%%%%%%%%%%%%%%%%%%%%%%%%%%%%%%%%%%%%%
% kda-regime-arity.tex
%
% Two-column conference-style paper. Built and verified with tectonic 0.16.9:
%   cd paper && tectonic -X compile kda-regime-arity.tex
%
% NOTE ON TOOLCHAIN: there is no conference .cls available in this environment,
% so the conference look is approximated with article + mathptmx + geometry.
% The bibliography is a manual thebibliography environment on purpose -- no
% bibtex/biber pass is required or run.
%%%%%%%%%%%%%%%%%%%%%%%%%%%%%%%%%%%%%%%%%%%%%%%%%%%%%%%%%%%%%%%%%%%%%%%%%%%%%%%

\documentclass[10pt,twocolumn]{article}

% ---- Fonts and page geometry -------------------------------------------------
% IMPORTANT: do NOT swap mathptmx for \usepackage{times} here. tectonic runs
% XeTeX, whose default encoding is TU; there is no TU-encoded `ptm', so
% \usepackage{times} silently falls back to Latin Modern and the document is
% not actually set in Times (verified with pdffonts: no Times glyph embedded).
\usepackage[T1]{fontenc}
\usepackage{mathptmx}
\usepackage[letterpaper,margin=0.75in]{geometry}
\setlength{\columnsep}{0.25in}
\usepackage{microtype}

% ---- Math and tables ---------------------------------------------------------
\usepackage{amsmath}
\usepackage{amssymb}
\usepackage{booktabs}
\usepackage{float}

% ---- URLs --------------------------------------------------------------------
\usepackage{xurl}

% ---- Author block ------------------------------------------------------------
\usepackage{authblk}
\renewcommand\Authfont{\normalsize}
\renewcommand\Affilfont{\small}
\setlength{\affilsep}{0.6em}

% ---- Float placement: keep tables near their callouts, not at the end --------
\renewcommand{\topfraction}{0.9}
\renewcommand{\bottomfraction}{0.8}
\renewcommand{\textfraction}{0.07}
\renewcommand{\floatpagefraction}{0.7}
\renewcommand{\dbltopfraction}{0.9}
\renewcommand{\dblfloatpagefraction}{0.7}
\setcounter{topnumber}{3}
\setcounter{bottomnumber}{2}
\setcounter{totalnumber}{4}
\setcounter{dbltopnumber}{3}

% ---- Slightly tighter conference-style section headings ----------------------
\makeatletter
\renewcommand\section{\@startsection{section}{1}{\z@}%
  {-2.0ex \@plus -0.5ex \@minus -0.2ex}{1.0ex \@plus 0.2ex}%
  {\normalfont\large\bfseries}}
\renewcommand\subsection{\@startsection{subsection}{2}{\z@}%
  {-1.6ex \@plus -0.5ex \@minus -0.2ex}{0.8ex \@plus 0.2ex}%
  {\normalfont\normalsize\bfseries}}
\makeatother

% ---- Table helper macros -----------------------------------------------------
\newcommand{\hd}[1]{\multicolumn{1}{c}{#1}}

\title{\bfseries Write Strength, Not Write Count: The $\beta$ Range Dominates\\
Householder Arity in Gated Delta-Rule Attention}

\author[1,2]{Eric Wu}
\affil[1]{Stanford University}
\affil[2]{Alpha AI Engineering}

\date{}

\begin{document}

%%%%%%%%%%%%%%%%%%%%%%%%%%%%%%%%%%%%%%%%%%%%%%%%%%%%%%%%%%%%%%%%%%%%%%%%%%%%%%%
% Full-width title + abstract block, conference style.
%%%%%%%%%%%%%%%%%%%%%%%%%%%%%%%%%%%%%%%%%%%%%%%%%%%%%%%%%%%%%%%%%%%%%%%%%%%%%%%
\makeatletter
\twocolumn[\begin{@twocolumnfalse}
\maketitle
\begin{center}
\begin{minipage}{0.86\textwidth}
\begin{abstract}
\noindent
Delta-rule linear attention with a channel-wise forget gate can be made more
expressive in two ways: by applying more Householder factors per token
(\emph{arity}, $R$), or by widening the range of the write-strength gate $\beta$
from the strict $(0,1)$ to the reflection range $(0,2)$. Recent work has
concentrated on the first. We report a small controlled factorial measurement
suggesting the second matters far more. In a fully paired $2\times2$ design
crossing $R\in\{1,2\}$ with $\beta$ regime, over two non-solvable group word
problems ($A_5$, $S_5$) and six seed bundles (48 cells, 3.52 GPU-hours), the
$\beta$ widening at fixed $R=1$ is large and significant in every position band
where both arms are above chance, while doubling $R$ under strict $\beta$ is
consistently positive but roughly two orders of magnitude smaller --- and the
$\beta$ widening is \emph{free}, adding zero parameters, whereas doubling $R$
costs $+40.3\%$. The $\beta\times R$ interaction is positive in five of eight
de-nested position bands. Because Kimi K3's KDA layer fixes
$\beta=\mathrm{Sigmoid}(\cdot)\in(0,1)$, a determinant argument shows K3-style
nodes cannot obtain a negative-determinant transition at any $R$; whether they
would benefit from $R>1$ in practice is a separate question our scale cannot
settle. We report four results against our own interest. Our pre-registered
mechanism --- that reflections help because $S_5$'s odd generator demands
$\det=-1$ --- is \emph{contradicted}, and the fallback explanation we offer is
contradicted too. Our $R$ arms are not parameter-matched. \emph{No} cell
survives Holm correction at $n=6$, and only one does after we correct a metric
artifact. And an asymmetric optimization failure (10 of 24 strict-$\beta$ runs
end with unconverged loss, versus 0 of 24 reflection runs) means we cannot
separate an expressivity effect from a trainability one. We report the effect,
its size, and the four reasons it is not yet established.
\end{abstract}
\end{minipage}
\end{center}
\vspace{1.5ex}
\end{@twocolumnfalse}]
\makeatother

%%%%%%%%%%%%%%%%%%%%%%%%%%%%%%%%%%%%%%%%%%%%%%%%%%%%%%%%%%%%%%%%%%%%%%%%%%%%%%%
\section{Introduction}
\label{sec:intro}

Linear-attention sequence models trade the quadratic cost of softmax attention
for a fixed-size recurrent state, and the delta rule
\cite{schlag2021fwp,yang2024deltanet} has become the dominant way to write into
that state: each token performs a rank-one \emph{erase} followed by a rank-one
\emph{write}. A well-documented consequence is limited state-tracking ability.
Two remedies have been proposed. \textsc{DeltaProduct}
\cite{siems2025deltaproduct} increases the number of Householder factors applied
per token, a quantity we call the \emph{arity} $R$. Separately, Grazzi et al.\
\cite{grazzi2024negative} showed that admitting \emph{negative} eigenvalues in
the state transition unlocks state tracking that is unreachable otherwise; in a
delta-rule node this is achieved by widening the write-strength gate $\beta$
from the strict range $(0,1)$ to the reflection range $(0,2)$, since $1-\beta$
is the non-unit eigenvalue of each factor.

These two levers are not independent, and the literature has largely pursued the
first. That matters practically, because arity is expensive --- each extra factor
needs its own key and gate projections --- while widening $\beta$ is a
one-character change to a sigmoid that adds no parameters at all. It also
matters now, because the largest deployed delta-rule attention variant, the Kimi
Delta Attention (KDA) layer of Kimi~K3 \cite{kimi2026k3}, specifies
$\beta_t = \mathrm{Sigmoid}(W_\beta x_t) \in (0,1)$ --- strictly inside the
non-reflective range.

We test this with a fully paired $2\times2$ factorial crossing arity
$R \in \{1,2\}$ with $\beta$ regime at otherwise identical geometry, on two
non-solvable group word problems, with six seed bundles per cell. The design is
deliberately narrow: it is built to estimate one interaction term rather than to
establish a state of the art. It is also small, and we have tried to be
scrupulous about what a 48-cell, $n=6$, $1.4$M-parameter probe can and cannot
support. Two things in particular govern how our results should be read, and we
put them here rather than in a late limitations section.

\paragraph{The metric is a prefix average, and we analyze around it.}
Our harness scores mean token accuracy over \emph{all} positions $1..L$ of each
evaluation sequence. Consequently the five nominal ``lengths'' are nested, and a
model that is correct below $L=40$ and at chance above it still produces a
smooth, gracefully-decaying curve that looks like partial extrapolation. For our
strict-$\beta$ arms this is exactly what happens: a zero-parameter model
asserting ``correct below 40, at chance above'' reproduces their entire length
curve to within $1.3$pp (Table~\ref{tab:dilution}). We therefore report
\emph{de-nested position bands} as our primary analysis
(Section~\ref{sec:denested}) and treat the prefix-averaged numbers as secondary.
De-nesting strengthens our $A_5$ result and \emph{eliminates} two of the seven
apparently significant $S_5$ cells, and we report both directions.

\paragraph{Both strict-$\beta$ arms are at chance beyond $L=40$.}
In every de-nested band past position 40, \texttt{R1} and \texttt{DP2-strict}
score within $2\times$ chance (Table~\ref{tab:bands}). This has a consequence we
must state plainly, because it defeats an argument we would otherwise have
liked to make: the strict-regime arity contrast cannot bound what capacity is
worth, since it is measured between two models that are both floored. The
parameter confound described in Section~\ref{sec:limitations} is therefore
\emph{open}, not bounded.

\paragraph{Contributions.}
\emph{(i)} A controlled factorial measurement of arity against $\beta$ range at
matched geometry, fully paired at the seed-bundle level so every arm within a
bundle is scored on a byte-identical evaluation bank (Section~\ref{sec:setup}).
\emph{(ii)} The finding that in every band where both arms are measurable, the
$\beta$ effect exceeds the arity effect by one to two orders of magnitude at
zero parameter cost (Section~\ref{sec:betavsr}), together with a positive
$\beta\times R$ interaction in five of eight de-nested bands --- of which one
survives Holm correction (Section~\ref{sec:denested}). \emph{(iii)} A
determinant observation about K3-style nodes, and a correction to a tempting but
wrong reading of K3's new bounded-decay parameterization
(Section~\ref{sec:k3}). \emph{(iv)} Two negative results against our own
hypotheses: the group-parity account of \emph{why} reflections help is
contradicted, and so is the headroom account we offered as a replacement
(Section~\ref{sec:parity}). \emph{(v)} An explicit statement of the
trainability confound that we cannot resolve, and of what would resolve it
(Section~\ref{sec:limitations}).

We claim no new operator and no new theory. The gated Householder operator we
measure is already implemented in public kernels and published at scale
\cite{li2026eda}; the determinant argument of Section~\ref{sec:det} is a special
case of a published result \cite{grazzi2024negative}, and weaker than it. What is
new is the controlled measurement. This is an empirical, partly negative result
and we frame it as one.

%%%%%%%%%%%%%%%%%%%%%%%%%%%%%%%%%%%%%%%%%%%%%%%%%%%%%%%%%%%%%%%%%%%%%%%%%%%%%%%
\section{Background}
\label{sec:background}

\subsection{The gated delta rule}

Let $S_t \in \mathbb{R}^{d_k \times d_v}$ be the recurrent state, $k_t$ the
(L2-normalized) key, $v_t$ the value, $\beta_t$ the write strength and
$\alpha_t \in (0,1)^{d_k}$ a channel-wise forget gate. The gated delta rule is
\begin{equation}
\label{eq:delta}
S_t = \bigl(I - \beta_t k_t k_t^{\top}\bigr)\,\mathrm{Diag}(\alpha_t)\,S_{t-1}
      + \beta_t k_t v_t^{\top},
\end{equation}
with readout $\tilde{o}_t = S_t^{\top} q_t$. Equation~\eqref{eq:delta} is
exactly the KDA recurrence of Kimi~K3, and specializes to gated DeltaNet
\cite{yang2024gated} and, with $\alpha$ scalar per head, to Mamba-2--style
gating \cite{dao2024mamba2}.

\subsection{Householder arity}

\textsc{DeltaProduct} \cite{siems2025deltaproduct} applies $R$ rank-one factors
per token, each with its own key and write strength:
\begin{equation}
\label{eq:arity}
S_t = \Bigl[\textstyle\prod_{i=1}^{R}
      \bigl(I - \beta_t^{(i)} k_t^{(i)} k_t^{(i)\top}\bigr)\Bigr]
      \mathrm{Diag}(\alpha_t)\, S_{t-1} + \text{writes}.
\end{equation}
Each factor is a true Householder reflection when $\beta = 2$.

\paragraph{Notation.} We write $R$ for the number of factors per token and call
it \emph{arity}. \textsc{DeltaProduct} writes $n_h$, denoting models
$\textsc{DeltaProduct}_{n_h}[\cdot]$ with the admissible eigenvalue range in
brackets, and does not use the word ``arity''; our $R$ is their $n_h$, and our
arms are their $n_h \in \{1,2\}$ (their $n_h{=}1$ is DeltaNet). Arity is not
free: the $R{=}2$ node needs a second set of key and gate projections.

\subsection{The determinant fact, and to whom it belongs}
\label{sec:det}

The per-token map acting on $S_{t-1}$ in Equation~\eqref{eq:arity} is
$A_t = \bigl[\prod_i (I - \beta^{(i)}_t k^{(i)}_t k^{(i)\top}_t)\bigr]
\mathrm{Diag}(\alpha_t)$. For a unit-norm key,
$\det(I - \beta kk^{\top}) = 1 - \beta$, so
\begin{equation}
\label{eq:det}
\det A_t = \Bigl[\textstyle\prod_{i=1}^{R}\bigl(1-\beta_t^{(i)}\bigr)\Bigr]
           \cdot \textstyle\prod_{j} \alpha_{t,j}.
\end{equation}
Since $\alpha_{t,j} > 0$ by construction, the \emph{sign} of $\det A_t$ is
governed entirely by the $\beta$ factors, and if every $\beta^{(i)} \in (0,1)$
then $\det A_t > 0$ \emph{for every} $R$: raising arity cannot recover a
negative determinant that the $\beta$ range has excluded.

\paragraph{What this does \emph{not} predict.}
We are careful here, because Equation~\eqref{eq:det} does not by itself predict
the interaction we go on to measure. At $R=2$ with $\beta\in(0,2)$, two negative
factors multiply to a \emph{positive} determinant, so $R{=}2$ reflection unlocks
no determinant sign that $R{=}1$ reflection cannot already reach. Read strictly,
Equation~\eqref{eq:det} predicts a main effect of $\beta$ and \emph{no}
$\beta\times R$ interaction at $R=2$. Our measured interaction is therefore not
explained by this argument, and we do not present it as confirmed theory. The
determinant argument licenses exactly one claim --- that strict $\beta$
forecloses negative determinants at every arity --- and we restrict ourselves to
that.

\paragraph{Attribution.}
Grazzi et al.\ \cite{grazzi2024negative} study $\mathcal{M}^n_k(\Omega)$, the
products of exactly $k$ generalized Householder matrices
$I - \beta_i v_i v_i^\top$ with $(1-\beta_i) \in \Omega$, $\|v_i\|=1$. Their
Proposition~1.3 states that any eigenvalue of $N \in \mathcal{M}^n_k((-1,1])$ is
either $1$ or has modulus below $1$, and that \emph{if in addition}
$N \in \mathcal{M}^n_k([0,1])$ \emph{and} $k \le 2$, then it lies in
$[0,1]\subset\mathbb{R}$. The $k \le 2$ hypothesis is essential --- it was added
in a published correction, and its failure for $k \ge 3$ is what
\textsc{DeltaProduct} exploits. They conclude such transitions are the identity
or non-orthogonal, hence neither reflections nor rotations, and that for
$k \le 2$ the model cannot learn parity by their Theorem~1. Their
Proposition~1.2 is the permutation statement proper: realizing a permutation
matrix requires the eigenvalue set $\{-1,1\}$. Equation~\eqref{eq:det} is a
determinant-level restatement, for the delta rule with normalized keys, of a
spectral fact they establish in far greater generality; it also says strictly
less, since at $R\ge3$ a positive determinant no longer implies non-negative
real eigenvalues. We claim no theoretical novelty for it whatsoever.

%%%%%%%%%%%%%%%%%%%%%%%%%%%%%%%%%%%%%%%%%%%%%%%%%%%%%%%%%%%%%%%%%%%%%%%%%%%%%%%
\section{The K3 KDA node}
\label{sec:k3}

Kimi~K3 \cite{kimi2026k3} is a mixture-of-experts model with $2.78$T total and
$104.2$B activated parameters, a 1M-token context, and native vision. Each block
places three KDA layers followed by one gated MLA layer --- a $3{:}1$ ratio
\emph{per block} (\S2.1) --- with one additional MLA layer at the end of the
backbone, giving 69 KDA and 24 MLA layers (Table~1, \S3.2) and a global ratio of
$2.875{:}1$. K3 uses no explicit positional encoding (\S3.4) and applies NoPE to
all MLA layers, so positional and recency information reaches the model only
through KDA's gating and decay (\S2.1.2). The KDA recurrence is
Equation~\eqref{eq:delta}, and its write strength is
\begin{equation}
\label{eq:k3beta}
\beta_t^{h} = \mathrm{Sigmoid}\bigl(W_\beta^{h} x_t\bigr) \in (0,1),
\end{equation}
a plain sigmoid, per head and per token, with no scale factor: we searched the
report and $\beta$ is never scaled, doubled, or rescaled anywhere, with the open
interval stated in two independent places (\S2.1.1). The remaining projections
are $q,k = \mathrm{L2Norm}(\mathrm{Swish}(\mathrm{ShortConv}(W_{q/k}x)))$,
$v = \mathrm{Swish}(\mathrm{ShortConv}(W_v x))$, and a decay logit
$z = W_\alpha^{\uparrow} W_\alpha^{\downarrow} x + b_\alpha^{h}$ from a low-rank
projection plus a per-head bias.

By Equation~\eqref{eq:det}, K3's KDA layers have $\det A_t > 0$ at every token
and head, and this would remain true at any $R$. That is a statement about what
strict $\beta$ forecloses, and it holds independently of our experiment. Whether
a K3-style node would \emph{in practice} benefit from $R>1$ is an empirical
question that our $1.4$M-parameter probe cannot settle; we return to this in
Section~\ref{sec:limitations}. We note the report contains no Householder,
arity, or \textsc{DeltaProduct} discussion, and no analysis of KDA's
state-tracking properties; the question we ask is not one K3 addresses.

\paragraph{A correction worth making explicitly.}
K3 introduces a \emph{bounded} decay parameterization (\S2.1.1, Eq.~5),
$g = g_{\min}\!\cdot\!\mathrm{Sigmoid}(e^{A_h} z)$ with $g_{\min}=-5$ fixed and
$A_h$ a learnable per-head log-scale initialized to $0$, so
$\alpha = \exp(g) \in (e^{g_{\min}}, 1)$. K3 reports this replaces the unbounded
negative-softplus mapping $g = -e^{A_h}\mathrm{Softplus}(z) \in (-\infty,0)$
used by Kimi~Linear \cite{kimi2025linear}. It is tempting to read this floor as
what forecloses reflections. It is not. Because $\alpha = \exp(g) > 0$ for any
finite $g$, $\det\mathrm{Diag}(\alpha) > 0$ under \emph{both} parameterizations,
so the floor cannot change the sign of $\det A_t$ at all. Determinant sign is
governed solely by $\beta$. What forecloses odd permutations in K3 is the
sigmoid-parameterized $\beta$ of Equation~\eqref{eq:k3beta}, which Kimi~Linear
already had.

K3's stated motivation for the floor is numerical, not representational. The
reciprocal cumulative decay $1/\Gamma$ used to rescale keys within a chunk can
grow without bound and overflow in finite precision; with $g_{\min}=-5$ the
cumulative log-decay over a 16-token tile is bounded, so $1/\Gamma$ stays within
BF16 range and \emph{all} causal tiles can use dense Tensor Core matrix
multiplications, eliminating the position-pair diagonal path (Fig.~3). No
modelling or quality justification is offered, and the report contains no
architecture ablations at all --- so this is an absence of evidence, not
evidence of absence. K3 does note the bounded form is related to lower-bounded
recurrence gates in prior work; we read that as a relatedness remark appended
after the numerical argument, since the paragraph's topic sentence and the
figure caption are both purely computational.

One consequence K3 does not discuss cuts against our own framing: the floor
$\alpha > e^{-5} \approx 6.7\times10^{-3}$ is a genuine modelling constraint,
since near-total single-step forgetting was representable under Kimi~Linear's
unbounded range and is not under K3's. Given that KDA is also K3's only source
of positional information, the decay range is load-bearing in a way the
numerical framing does not convey. We raise it rather than let it be raised
against us.

%%%%%%%%%%%%%%%%%%%%%%%%%%%%%%%%%%%%%%%%%%%%%%%%%%%%%%%%%%%%%%%%%%%%%%%%%%%%%%%
\section{Experimental setup}
\label{sec:setup}

\paragraph{Arms.}
We cross arity $R \in \{1,2\}$ with $\beta$ regime, giving four arms
(Table~\ref{tab:arms}). All use the same gated Householder mixer with the same
Triton backend, 4000 steps, batch 64, and no feed-forward blocks. The $\beta$
regime is a single scale on the sigmoid, $\beta = s\cdot\mathrm{Sigmoid}(z)$
with $s\in\{1,2\}$; it adds no parameters, and the two regimes have
byte-identical parameter ledgers. Post hoc, the strict arms reach
$\beta_{\max}=1.0000$ and the reflection arms $\beta_{\max}=2.0000$, with no
record violating its bound.

\paragraph{Tasks.}
Two group word problems: given a sequence of generators of a finite group,
predict the composed element. We use $A_5$ (order 60, non-solvable, \emph{both}
generators even) and $S_5$ (order 120, non-solvable, transposition generator
\emph{odd}). Chance accuracy is $1/60 = 1.67\%$ and $1/120 = 0.83\%$
respectively. The parity difference motivated the prediction
Section~\ref{sec:parity} reports as falsified.

\paragraph{The metric, stated explicitly.}
Accuracy is the fraction of correct predictions averaged over \emph{all} $L$
positions of 64 evaluation sequences ($N = 64L$ scored tokens; we verified all
240 reported values are integer multiples of $1/64L$). Training words have
length $3$--$40$. Because the metric averages positions $1..L$, the nominal
lengths are \emph{nested}: $\mathrm{acc}(512)$ includes positions $\le 40$ that
are in-distribution. We therefore define \emph{de-nested bands}
\begin{equation}
\label{eq:denest}
\mathrm{band}(L_{i-1},L_i) =
\frac{L_i\,\mathrm{acc}(L_i) - L_{i-1}\,\mathrm{acc}(L_{i-1})}{L_i - L_{i-1}},
\end{equation}
the mean accuracy on positions in $(L_{i-1}, L_i]$ alone, and use these as the
primary analysis. Only bands past position 40 are extrapolation.

\paragraph{Pairing, and why $n=6$ helps.}
Six seed bundles, $1101$--$1106$, each fixing four independent seed streams
(initialization, data, task, evaluation), pairwise distinct in every record.
Within a bundle all four arms are scored on the same evaluation bank: we
confirmed one distinct \texttt{eval\_bank\_sha256} per (task, bundle) shared by
all four arms, and twelve mutually distinct hashes across the twelve groups.
Every contrast is a within-bundle difference cancelling evaluation-sampling
noise identically in all arms. There were $0$ train/eval collisions.

\paragraph{Statistics.}
For each cell we compute per-bundle differences and aggregate those. The
interaction is differenced \emph{per bundle before averaging},
$i_b = d_b^{\text{refl}} - d_b^{\text{strict}}$, preserving pairing; computing
it from arm means with an unpaired standard error would be wrong. We report
mean, $\mathrm{se}=\mathrm{sd}/\sqrt6$ (sd with $\mathrm{ddof}=1$),
$t=\mathrm{mean}/\mathrm{se}$, and the one-sided 95\% lower bound
$\mathrm{L95}=\mathrm{mean}-t_{0.95,5}\,\mathrm{se}$, $t_{0.95,5}=2.015$. We
also report Holm-corrected outcomes across the band family, since we test many
cells.

\paragraph{Compute and provenance.}
All 48 cells ran on one NVIDIA A10G, $48/48$ completing, for $3.52$ GPU-hours
(cells $228$--$305$\,s). All 48 carry one probe-source revision, $0$ records
were rejected, $0$ non-finite loss steps occurred.

%%%%%%%%%%%%%%%%%%%%%%%%%%%%%%%%%%%%%%%%%%%%%%%%%%%%%%%%%%%%%%%%%%%%%%%%%%%%%%%
% TABLE: arms
%%%%%%%%%%%%%%%%%%%%%%%%%%%%%%%%%%%%%%%%%%%%%%%%%%%%%%%%%%%%%%%%%%%%%%%%%%%%%%%
\begin{table}[t]
\centering
\footnotesize
\caption{The four arms. ``Mixer params'' excludes both the token embedding and
the output head; it is task-invariant, whereas total parameters differ by task
because the head does ($+39.7\%$ on $A_5$, $+39.1\%$ on $S_5$ from $R{=}1$ to
$R{=}2$). The $\beta$ regime is parameter-free. $\bar\beta$ is the mean realized
$\beta$ pooled over both tasks and all six bundles.}
\label{tab:arms}
\begin{tabular}{llrrr}
\toprule
Arm & $\beta$ range & \hd{$R$} & \hd{Mixer params} & \hd{$\bar\beta$} \\
\midrule
\texttt{R1}          & $(0,1)$ & 1 & 998{,}092   & 0.899 \\
\texttt{DP2-strict}  & $(0,1)$ & 2 & 1{,}400{,}524 & 0.779 \\
\texttt{R1-refl}     & $(0,2)$ & 1 & 998{,}092   & 1.800 \\
\texttt{Reflection}  & $(0,2)$ & 2 & 1{,}400{,}524 & 1.658 \\
\bottomrule
\end{tabular}
\end{table}

%%%%%%%%%%%%%%%%%%%%%%%%%%%%%%%%%%%%%%%%%%%%%%%%%%%%%%%%%%%%%%%%%%%%%%%%%%%%%%%
% TABLE 1 -- full width: accuracy, prefix and de-nested
%%%%%%%%%%%%%%%%%%%%%%%%%%%%%%%%%%%%%%%%%%%%%%%%%%%%%%%%%%%%%%%%%%%%%%%%%%%%%%%
\begin{table*}[t]
\centering
\small
\caption{Mean accuracy (\%), $n=6$ bundles. \textbf{Left:} the prefix-averaged
metric our harness reports, over positions $1..L$. \textbf{Right:} the same data
de-nested into position bands via Equation~\eqref{eq:denest}, which is what
actually measures behaviour at new lengths. Chance is $1.67\%$ ($A_5$) and
$0.83\%$ ($S_5$); \emph{italic} marks band values within $2\times$ chance. Both
strict arms are at chance in every extrapolation band --- the prefix metric
disguises this as graceful decay.}
\label{tab:bands}
\begin{tabular}{llrrrrr@{\hspace{2em}}rrrrr}
\toprule
& & \multicolumn{5}{c}{Prefix-averaged, positions $1..L$}
  & \multicolumn{5}{c}{De-nested bands, positions in $(L_{i-1},L_i]$} \\
\cmidrule(lr){3-7}\cmidrule(lr){8-12}
Arm & Task & \hd{40} & \hd{64} & \hd{128} & \hd{256} & \hd{512}
        & \hd{1--40} & \hd{41--64} & \hd{65--128} & \hd{129--256} & \hd{257--512} \\
\midrule
\texttt{R1}         & $A_5$ & 79.19 & 49.96 & 25.74 & 13.83 & 7.82
                    & 79.19 & \emph{1.24} & \emph{1.53} & \emph{1.91} & \emph{1.81} \\
\texttt{DP2-strict} & $A_5$ & 82.49 & 53.45 & 27.75 & 14.81 & 8.17
                    & 82.49 & 5.05 & \emph{2.04} & \emph{1.87} & \emph{1.53} \\
\texttt{R1-refl}    & $A_5$ & 99.99 & 99.70 & 82.83 & 46.47 & 23.63
                    & 99.99 & 99.22 & 65.97 & 10.11 & \emph{0.78} \\
\texttt{Reflection} & $A_5$ & 100.00 & 100.00 & 98.31 & 82.73 & 54.99
                    & 100.00 & 100.00 & 96.62 & 67.15 & 27.25 \\
\cmidrule(lr){1-12}
\texttt{R1}         & $S_5$ & 74.06 & 46.98 & 23.65 & 12.51 & 6.63
                    & 74.06 & \emph{1.86} & \emph{0.31} & \emph{1.37} & \emph{0.76} \\
\texttt{DP2-strict} & $S_5$ & 76.73 & 48.90 & 24.47 & 12.89 & 6.84
                    & 76.73 & \emph{2.53} & \emph{0.04} & \emph{1.31} & \emph{0.79} \\
\texttt{R1-refl}    & $S_5$ & 97.07 & 72.86 & 37.30 & 19.58 & 9.99
                    & 97.07 & 32.50 & \emph{1.75} & \emph{1.85} & \emph{0.41} \\
\texttt{Reflection} & $S_5$ & 98.09 & 88.45 & 56.04 & 28.53 & 14.64
                    & 98.09 & 72.40 & 23.62 & \emph{1.02} & \emph{0.75} \\
\bottomrule
\end{tabular}
\end{table*}

%%%%%%%%%%%%%%%%%%%%%%%%%%%%%%%%%%%%%%%%%%%%%%%%%%%%%%%%%%%%%%%%%%%%%%%%%%%%%%%
% TABLE: dilution fit
%%%%%%%%%%%%%%%%%%%%%%%%%%%%%%%%%%%%%%%%%%%%%%%%%%%%%%%%%%%%%%%%%%%%%%%%%%%%%%%
\begin{table}[t]
\centering
\footnotesize
\caption{Why we de-nest. Maximum absolute error, over
$L\in\{64,128,256,512\}$, of the \emph{zero-free-parameter} prediction
$\mathrm{acc}(L) = \frac{40}{L}\mathrm{acc}(40)
+ \frac{L-40}{L}\cdot\text{chance}$, i.e.\ ``correct below 40, at chance
above''. The strict arms are fit to within $1.3$pp: their apparent length curve
carries no extrapolation signal. The reflection arms deviate by $12$--$66$pp, so
they genuinely do extrapolate.}
\label{tab:dilution}
\begin{tabular}{lrr}
\toprule
Arm & \hd{$A_5$ max err} & \hd{$S_5$ max err} \\
\midrule
\texttt{R1}          & 0.16 & 0.38 \\
\texttt{DP2-strict}  & 1.27 & 0.64 \\
\texttt{R1-refl}     & 50.44 & 11.87 \\
\texttt{Reflection}  & 65.92 & 26.84 \\
\bottomrule
\end{tabular}
\end{table}

%%%%%%%%%%%%%%%%%%%%%%%%%%%%%%%%%%%%%%%%%%%%%%%%%%%%%%%%%%%%%%%%%%%%%%%%%%%%%%%
% TABLE 2 -- full width: interaction, both metrics
%%%%%%%%%%%%%%%%%%%%%%%%%%%%%%%%%%%%%%%%%%%%%%%%%%%%%%%%%%%%%%%%%%%%%%%%%%%%%%%
\begin{table*}[t]
\centering
\small
\caption{The $\beta\times R$ interaction, differenced per bundle before
averaging. \textbf{Left:} on the prefix-averaged metric. \textbf{Right:} on
de-nested position bands, our primary analysis. L95 is the one-sided 95\% lower
bound ($t_{0.95,5}=2.015$); \textbf{bold} marks L95 $>0$ uncorrected. Seven of
ten prefix cells and five of eight bands are positive at uncorrected L95. Under
Holm correction across the family, \textbf{zero} of the ten prefix cells and
\textbf{one} of the eight bands survive (marked $\dagger$): with $df=5$,
$p<0.00625$ needs $|t|>4.03$ and our maximum band $t$ is $4.26$. Note $S_5$ at
$129$--$256$ and $257$--$512$ \emph{reverse sign} once de-nested --- they were
significant only by inheriting the $41$--$64$ signal through prefix averaging.}
\label{tab:paired}
\begin{tabular}{lrrrr@{\hspace{2.2em}}lrrrr}
\toprule
\multicolumn{5}{c}{Prefix-averaged, positions $1..L$}
& \multicolumn{5}{c}{De-nested bands} \\
\cmidrule(lr){1-5}\cmidrule(lr){6-10}
Task, $L$ & \hd{Inter.} & \hd{se} & \hd{$t$} & \hd{L95}
& Task, band & \hd{Inter.} & \hd{se} & \hd{$t$} & \hd{L95} \\
\midrule
$A_5$, 40  & $-3.30$ & 4.25 & $-0.78$ & $-11.87$
  & $A_5$, 41--64   & $-3.03$ & 1.53 & $-1.97$ & $-6.12$ \\
$A_5$, 64  & $-3.20$ & 3.09 & $-1.03$ & $-9.43$
  & $A_5$, 65--128  & $\mathbf{+30.14}$ & 9.34 & $3.23$ & $\mathbf{+11.32}$ \\
$A_5$, 128 & $\mathbf{+13.47}$ & 4.74 & $2.84$ & $\mathbf{+3.93}$
  & $A_5$, 129--256 & $\mathbf{+57.09}$ & 16.08 & $3.55$ & $\mathbf{+24.69}$ \\
$A_5$, 256 & $\mathbf{+35.28}$ & 9.38 & $3.76$ & $\mathbf{+16.37}$
  & $A_5$, 257--512 & $\mathbf{+26.74}$ & 12.56 & $2.13$ & $\mathbf{+1.43}$ \\
$A_5$, 512 & $\mathbf{+31.01}$ & 10.70 & $2.90$ & $\mathbf{+9.45}$
  & $S_5$, 41--64   & $\mathbf{+39.23}^{\dagger}$ & 9.21 & $4.26$ & $\mathbf{+20.66}$ \\
$S_5$, 40  & $-1.65$ & 2.42 & $-0.68$ & $-6.52$
  & $S_5$, 65--128  & $\mathbf{+22.14}$ & 8.80 & $2.52$ & $\mathbf{+4.40}$ \\
$S_5$, 64  & $\mathbf{+13.68}$ & 4.23 & $3.23$ & $\mathbf{+5.15}$
  & $S_5$, 129--256 & $-0.77$ & 0.79 & $-0.98$ & $-2.35$ \\
$S_5$, 128 & $\mathbf{+17.91}$ & 6.33 & $2.83$ & $\mathbf{+5.15}$
  & $S_5$, 257--512 & $+0.31$ & 0.24 & $1.31$ & $-0.17$ \\
$S_5$, 256 & $\mathbf{+8.57}$ & 3.19 & $2.69$ & $\mathbf{+2.15}$
  & & & & & \\
$S_5$, 512 & $\mathbf{+4.44}$ & 1.68 & $2.64$ & $\mathbf{+1.05}$
  & & & & & \\
\bottomrule
\end{tabular}
\end{table*}

%%%%%%%%%%%%%%%%%%%%%%%%%%%%%%%%%%%%%%%%%%%%%%%%%%%%%%%%%%%%%%%%%%%%%%%%%%%%%%%
\section{Results}
\label{sec:results}

\subsection{Accuracy, and what the prefix metric hides}
\label{sec:denested}

Table~\ref{tab:bands} gives both views. On the prefix-averaged metric the
dominant feature is the $\beta$ contrast: at $L=128$ on $A_5$, moving from
strict to reflection $\beta$ at fixed $R=1$ takes accuracy from $25.74\%$ to
$82.83\%$, while doubling $R$ at fixed strict $\beta$ takes it from $25.74\%$ to
$27.75\%$.

De-nesting sharpens this and also disciplines it. In every extrapolation band,
both strict arms sit within $2\times$ chance --- \texttt{R1} on $A_5$ scores
$1.24$--$1.91\%$ against $1.67\%$ chance. The reflection arms are genuinely
above chance for longer: \texttt{Reflection} on $A_5$ reaches $96.62\%$ on
positions $65$--$128$ and $27.25\%$ on $257$--$512$. Table~\ref{tab:dilution}
makes the point quantitatively: a zero-parameter ``correct below 40, at chance
above'' model reproduces both strict arms' entire length curves to within
$1.3$pp, and fails on the reflection arms by $12$--$66$pp. The honest reading is
that \emph{the strict-$\beta$ arms do not extrapolate past the training range at
all}, and the reflection arms partly do.

This also corrects an interpretation we would otherwise have offered. The
strict-regime arity effect on the prefix metric decays monotonically with length
($A_5$: $+3.31, +3.50, +2.01, +0.98, +0.35$pp), which invites a story about
extra factors mattering less as the task hardens. It is dilution: the pure
prediction $(40/L)\times\text{effect}(40)$ gives $S_5$
$+2.67, +1.67, +0.83, +0.42, +0.21$ against observed
$+2.67, +1.92, +0.82, +0.38, +0.21$. That curve is one number divided by $L$,
and we withdraw any interpretation of its shape.

\subsection{The interaction}

Table~\ref{tab:paired} gives the interaction under both metrics. On de-nested
bands it is positive at uncorrected L95 in five of eight bands, from $+22.14$ to
$+57.09$pp. De-nesting \emph{strengthens} $A_5$ ($+13.47 \to +30.14$ at
$65$--$128$; $+35.28 \to +57.09$ at $129$--$256$) and \emph{eliminates} the two
long-length $S_5$ cells, which reverse sign ($+8.57 \to -0.77$;
$+4.44 \to +0.31$). Those two cells were significant only because prefix
averaging carried the $41$--$64$ signal forward; their per-bundle interaction
values correlate at $\rho \ge 0.99$ with the $L=64$ cell. We do not count them.

We must be equally direct about multiplicity. Across the ten prefix cells,
\emph{no} cell survives Holm or Bonferroni correction --- Holm stops at the
first comparison ($p=0.00658$ against a threshold of $0.005$). This is
structural, not unlucky: with $df=5$, clearing $\alpha=0.005$ requires
$|t|>4.03$, and no prefix cell reaches it. On the eight de-nested bands, exactly
one survives ($S_5$, $41$--$64$, $t=4.26$, $p=0.00401$). So the defensible
statement is: one band-level interaction is established at family-wise
$\alpha=0.05$, and four more are positive and large but individually
uncorrected. The count-based framing ``$k$ of $n$ cells significant'' is
especially weak here because the cells are not independent (mean within-task
cross-length $r=0.53$, several at $0.99$--$1.00$).

\subsection{$\beta$ dominates arity, and it is free}
\label{sec:betavsr}

Table~\ref{tab:cost} compares the two levers in the bands where both arms are
measurable. On $A_5$ positions $41$--$64$, widening $\beta$ at fixed $R=1$ buys
$+97.98$pp ($t=87.75$) against arity's $+3.81$pp ($t=2.40$), a ratio of
$25.7\times$; on $65$--$128$, $+64.44$ against $+0.52$, or $124.7\times$; on
$S_5$ $41$--$64$, $+30.64$ against $+0.67$, or $45.5\times$. Beyond those bands
the comparison stops being meaningful because all four arms approach chance, and
we report no ratio there. On the prefix metric the corresponding ratios run
$6.3\times$ to $44.9\times$ across all ten cells, but we now regard those as
contaminated and give them only for continuity.

The cost asymmetry is the part we consider most robust, because it does not
depend on the effect size at all: the $\beta$ widening adds \emph{zero}
parameters --- both arms have $998{,}092$ mixer parameters and byte-identical
ledgers --- whereas doubling arity costs $+40.3\%$ of mixer parameters
($+39.7\%$ / $+39.1\%$ of total, by task). Whatever the true effect sizes, one
intervention is free and the other is not.

\paragraph{We do not claim arity does nothing.}
An earlier framing of this work said arity's effect was ``not distinguishable
from zero''. That is true cell-by-cell and misleading overall. The strict-regime
arity effect is positive in $10$ of $10$ prefix cells (one-sided sign test
$p=0.00098$, more significant than any single interaction cell we report), and
pooling all 60 per-bundle differences gives $+1.61$pp, $\mathrm{se}=0.61$,
$t=2.66$ --- though those 60 values are nested and not independent, so the
pooled test overstates its own precision. Our design also had little power: the
minimum detectable effect at 80\% power is $7.4$--$12.5$pp at short lengths, so
a real $+3.3$pp arity effect would have been invisible to us. The accurate
statement is that \emph{arity's benefit under strict $\beta$ is consistently
positive, one to two orders of magnitude smaller than the $\beta$ effect, and
below this design's per-cell resolution} --- not that it is absent.

%%%%%%%%%%%%%%%%%%%%%%%%%%%%%%%%%%%%%%%%%%%%%%%%%%%%%%%%%%%%%%%%%%%%%%%%%%%%%%%
% TABLE 3 -- single column
%%%%%%%%%%%%%%%%%%%%%%%%%%%%%%%%%%%%%%%%%%%%%%%%%%%%%%%%%%%%%%%%%%%%%%%%%%%%%%%
\begin{table}[t]
\centering
\footnotesize
\caption{$\beta$ range versus arity, on de-nested bands, restricted to the three
extrapolation bands where the compared arms are above chance. Both interventions
start from the same $R{=}1$ strict baseline; $\Delta$ is a paired mean
difference in points. The $\beta$ intervention is parameter-free; arity costs
$+40.3\%$ mixer parameters. Bands past these are omitted because all arms
approach chance, making the ratio meaningless.}
\label{tab:cost}
\begin{tabular}{llrrr}
\toprule
Task, band & Lever & \hd{$\Delta$} & \hd{$t$} & \hd{Ratio} \\
\midrule
$A_5$, 41--64  & $\beta$ & $+97.98$ & 87.75 & \\
$A_5$, 41--64  & arity   & $+3.81$  & 2.40  & $25.7\times$ \\
\cmidrule(lr){1-5}
$A_5$, 65--128 & $\beta$ & $+64.44$ & 7.55  & \\
$A_5$, 65--128 & arity   & $+0.52$  & 1.36  & $124.7\times$ \\
\cmidrule(lr){1-5}
$S_5$, 41--64  & $\beta$ & $+30.64$ & 6.05  & \\
$S_5$, 41--64  & arity   & $+0.67$  & 1.02  & $45.5\times$ \\
\bottomrule
\end{tabular}
\end{table}

%%%%%%%%%%%%%%%%%%%%%%%%%%%%%%%%%%%%%%%%%%%%%%%%%%%%%%%%%%%%%%%%%%%%%%%%%%%%%%%
\section{Two failed explanations}
\label{sec:parity}

We stated a mechanistic prediction before running the sweep. It failed. The
replacement we then offered also fails against the same data. We report both.

\paragraph{The parity prediction, and its falsification.}
Equation~\eqref{eq:det} says reflection $\beta$ unlocks negative determinants,
i.e.\ odd permutations. Our tasks differ exactly in whether odd permutations are
demanded: $S_5$'s transposition generator is odd, so composing a word requires
elements of determinant $-1$; in $A_5$ both generators are even. We predicted a
\emph{large} interaction on $S_5$ and a \emph{small} one on $A_5$. We observe the
opposite in the bands that matter: de-nested, the interaction is $+30.14$ and
$+57.09$pp on $A_5$ at $65$--$128$ and $129$--$256$, against $-0.77$ and
$+0.31$pp on $S_5$ in the same bands. Only the $41$--$64$ band favours $S_5$.
The reachability-of-reflections story does not explain our result. As
Section~\ref{sec:det} notes, it never should have: Equation~\eqref{eq:det}
predicts no interaction at $R=2$ in the first place.

\paragraph{The headroom explanation, also contradicted.}
We then proposed that $\beta>1$ helps by making each write more expressive
regardless of parity, with $A_5$ (order 60) simply easier than $S_5$ (order 120)
so its reflection arms retain headroom where $S_5$'s have collapsed. This
predicts the interaction should track available headroom. It does not. Across
the ten prefix cells, the correlation between interaction and headroom
($100 - \mathrm{acc}$ of \texttt{R1-refl}) is only $r=+0.53$; within $A_5$ it
fits well ($r=+0.93$) but within $S_5$ there is no relationship ($r=+0.28$,
$\rho=+0.10$), and using the \texttt{Reflection} arm's de-nested band accuracy
as the regressor gives the \emph{wrong sign} ($r=-0.33$). The specific claim we
had drafted --- that the $S_5$ reflection arms at $19.58\%$ ``have little room
left'' --- is also indefensible: $19.58\%$ against a $0.83\%$ chance floor is
$\sim$80pp of headroom, more than $A_5$ at $L=256$ ($46.47\%$), where the
interaction is four times larger.

\paragraph{Where that leaves the mechanism.}
We have a measured interaction, a falsified pre-registered mechanism, a
falsified fallback, and a determinant argument that predicts no interaction at
$R=2$. We therefore offer \emph{no} mechanism. A reader should treat the
interaction as an unexplained empirical regularity, and should weigh it against
the trainability account in Section~\ref{sec:limitations}, which we cannot rule
out. Distinguishing the candidates needs a solvable-group control, a
difficulty-matched non-solvable pair, and an $R=3$ arm (where
Equation~\eqref{eq:det} does predict something, since the $k\le2$ hypothesis of
Grazzi Prop.~1.3 fails there). We ran none of these.

%%%%%%%%%%%%%%%%%%%%%%%%%%%%%%%%%%%%%%%%%%%%%%%%%%%%%%%%%%%%%%%%%%%%%%%%%%%%%%%
\section{Limitations}
\label{sec:limitations}

\subsection{The trainability confound, which we cannot resolve}

This is the objection we consider most serious. Ten of 48 runs end with a final
logged loss above $0.1$, and \emph{all ten are strict-$\beta$} (\texttt{R1}
6/12, \texttt{DP2-strict} 4/12; both reflection arms 0/12). A perfectly
asymmetric stability failure between the arms being compared is the signature of
an optimization difference, not only an expressivity one. Within the 24 strict
runs, log final loss correlates with accuracy at $r=-0.43$ ($L=40$), and spiked
runs are worse at every length.

The confound is load-bearing. Restricting the interaction to bundles where all
four arms ended below $0.1$ ($A_5$: 1102--1105, $n=4$; $S_5$: 1104--1105,
$n=2$), \emph{zero} of ten cells remain significant. The $A_5$ point estimates
survive largely intact ($+13.56$, $+28.52$, $+26.93$ at $L=128/256/512$ against
$+13.47$, $+35.28$, $+31.01$), which is the best we can say for the effect; the
significance does not. Traces make the problem concrete: \texttt{R1} on $S_5$
bundle 1106 ends at $0.85$ having reached $10^{-4}$ at step 3000, and bundle
1103 oscillates between $\sim\!0$ and $1.27$ across the last 1500 steps. Those
are not converged models. Since $\beta=2$ doubles the effective write-gate
gradient scale, the reflection arms' perfect stability record may itself be an
optimization-geometry effect.

We therefore cannot separate ``reflection $\beta$ is more expressive'' from
``reflection $\beta$ is easier to optimize in 4000 steps''. Our loss traces are
logged only every 500 steps, too coarse to characterize the spikes, and we
retained no intermediate checkpoints, so we cannot re-evaluate from a
pre-spike model. Resolving this needs a learning-rate sweep per arm, longer
training, or best-checkpoint evaluation. Until then, this paper measures
\emph{the joint effect of expressivity and trainability} of the $\beta$ range,
and the reader should read every result above with that substitution.

\subsection{The $R$ arms are not parameter-matched, and the confound is open}

The $R{=}1$ arms have $998{,}092$ mixer parameters and the $R{=}2$ arms
$1{,}400{,}524$, $+40.3\%$. Any raw ``$R{=}2$ beats $R{=}1$'' main effect is
confounded with capacity and we claim none.

We had intended to argue that the strict-regime arity effect ($\le +3.50$pp)
upper-bounds what that capacity is worth. We withdraw that argument. As
Table~\ref{tab:bands} shows, both strict arms are at chance in every
extrapolation band, and a capacity contrast measured between two floored models
is guaranteed to read $\approx0$ regardless of what capacity is worth. The bound
is uninformative rather than reassuring, and using it would have been circular.
Two further problems: the \texttt{DP2-strict} vs \texttt{R1} contrast varies
arity and capacity together, so if arity were mildly harmful under strict
$\beta$ while capacity helped, the observed difference would bound neither
component; and the design's minimum detectable effect is $7.4$--$12.5$pp, so
capacity could be worth up to $\sim\!+11.9$pp within our confidence bounds.

It remains true that capacity differs by the same $402{,}432$ parameters in both
regimes, so to first order it cancels from a difference of differences. But that
argument assumes the \emph{value} of capacity is regime-independent, and we never
varied capacity within the reflection regime, so we cannot check it. The fix is
specified: the project defines an \texttt{R1-P} arm feed-forward-width-matched
to \texttt{DP2-strict} at $1{,}399{,}756$ parameters (a $0.055\%$ match), which
was \emph{not} run. Running it plus a reflection-regime counterpart --- roughly
24 cells at a cost comparable to this sweep --- would isolate capacity within
each regime. That control is necessary, not optional, and its absence is what
prevents us from making an arity claim.

\subsection{$n=6$, multiplicity, and influential bundles}

Six bundles makes paired contrasts informative but imprecise: interaction
standard errors run $1.68$--$16.08$pp, and the \texttt{Reflection} arm's
between-bundle spread on $A_5$ is large (sd $20.19$pp at $L=256$). As
Section~\ref{sec:denested} states, no prefix cell and only one band survives
Holm correction, and with $df=5$ none could have cleared Bonferroni even in
principle. Leave-one-bundle-out shows the two long-length $S_5$ prefix cells
lose significance if either of two bundles is dropped --- consistent with their
sign reversal under de-nesting. The $A_5$ result is more robust: dropping the
most extreme bundle (1103) takes $L=256$ from $+35.28$ to $+29.18$
(L95 $+10.56$) and $L=512$ from $+31.01$ to $+22.26$ (L95 $+6.18$), both still
positive. We report L95 values, not point estimates, as the defensible claim.

\subsection{Ceiling and floor censoring}

Of 240 reported values, 26 are exactly $100.00\%$ and 35 exceed $99\%$; of the
de-nested band values, $107$ of $240$ lie within $2\times$ chance and 10
band-cells fall strictly below it. At $L=40$ the reflection arms score $99.99\%$
and $100.00\%$ on $A_5$, the latter with zero variance across all six bundles,
so the $L=40$ interaction is mechanically bounded and its negative sign
($-3.30$pp) is not interpretable. As a robustness check we recomputed the
interaction on Haldane--Anscombe-corrected log-odds: sign and significance
broadly survive (9 of 10 prefix cells positive at uncorrected L95), but the
largest interaction moves from $L=256$ to $L=128$ on $A_5$, and $S_5$ at $L=40$
becomes significantly \emph{positive}. So the raw-scale magnitudes are partly an
artifact of pp differences being maximized near $50\%$, and our reading of the
$L=40$ cells as uninformative is right on $A_5$ but wrong on $S_5$.

\subsection{Scale, scope, and the K3 claim}

Our models have $1.0$--$1.4$M mixer parameters, no feed-forward blocks, and
train for 4000 steps on synthetic group word problems. Kimi~K3 has $2.78$T total
parameters. \textbf{Nothing in this paper is direct evidence about the behaviour
of a frontier-scale model.} We separate the two K3 statements accordingly. The
determinant observation --- that strict $\beta$ forecloses negative-determinant
transitions at every arity --- follows from Equation~\eqref{eq:det} alone, is
regime-agnostic, and does not depend on this experiment; we consider it sound.
The empirical prediction that K3-style nodes would not \emph{benefit} from
$R>1$ does depend on this experiment, and given the four caveats above we
present it as a hypothesis motivating a scaled replication, not as a finding.
Group word problems are also a deliberately adversarial probe chosen because
they discriminate; language modelling may weight these mechanisms differently.

We also flag a tension with the prior work we rely on.
\textsc{DeltaProduct} reports that its models ``fail to learn even the training
context-length when restricted to the standard eigenvalue range $[0,1]$,
regardless of the number of Householder transformations $n_h$''. Our strict-$\beta$
arms do learn the training range ($74$--$82\%$ at $L=40$). Either the tasks
differ materially or our strict arms are not the regime they describe; we cannot
say which.

\subsection{Other limitations}

We measure $R\in\{1,2\}$ only, so we cannot speak to $R\ge3$ --- where, as noted,
the theory actually predicts something different. All arms share one mixer
implementation and Triton backend, so a kernel bug interacting with the $\beta$
scale would not be caught by our design. We report a single learning-rate
schedule per arm with no per-arm tuning, which is precisely what the trainability
confound requires.

%%%%%%%%%%%%%%%%%%%%%%%%%%%%%%%%%%%%%%%%%%%%%%%%%%%%%%%%%%%%%%%%%%%%%%%%%%%%%%%
\section{Related work}
\label{sec:related}

\paragraph{Delta-rule linear attention.}
The delta rule as a fast-weight update goes back to Schlag et al.\
\cite{schlag2021fwp}; DeltaNet \cite{yang2024deltanet} made it trainable at
scale with a chunkwise-parallel form, and gated DeltaNet \cite{yang2024gated}
added the channel-wise forget gate that makes Equation~\eqref{eq:delta} the
modern default. Mamba-2 \cite{dao2024mamba2} supplies the structured-state-space
view of the same gating. Kimi~Linear \cite{kimi2025linear} introduced KDA, and
Kimi~K3 \cite{kimi2026k3} deploys it as the majority attention layer at frontier
scale.

\paragraph{Arity and negative eigenvalues are one research programme.}
The two lines of work our paper sits between share most of their authors.
Grazzi et al.\ \cite{grazzi2024negative} established that admitting negative
eigenvalues unlocks state tracking otherwise out of reach --- the result our
Section~\ref{sec:det} specializes and which subsumes our determinant
observation. \textsc{DeltaProduct} \cite{siems2025deltaproduct} then raised
arity, and its \S5.2 reports the finding that most directly prefigures ours:
``Unexpectedly, $S_4$ and $A_5$ can extrapolate robustly using only $n_h = 2$
despite the theorem suggesting 3 and 4, respectively.'' Two details matter for
us. They attribute the efficiency to $S_4$ and $A_5$ being isomorphic to
subgroups of $SO(3,\mathbb{R})$, and they observe a head that learns $\beta=2$
on \emph{both} factors --- two true reflections composing to a rotation. And
their \S5.2 uses the extended eigenvalue range throughout, reporting the failure
under $[0,1]$ quoted in Section~\ref{sec:limitations}. That sentence is the
closest prior statement to our result. Our contribution is to convert it from a
training-time observation reported in passing into a paired factorial estimate
at length extrapolation --- with the caveats above. Our result is a refinement of
their programme, not a challenge to it.

\paragraph{The operator is not new, and we are precise about what is prior.}
The per-channel-gated delta-rule operator we measure is computed by public
kernels for generalized (diagonal-plus-low-rank) delta-rule updates, and is
established at scale: EDA \cite{li2026eda} trains dense 2.5B and MoE
25B-A2.8B models on it --- though it appears there as the \emph{baseline},
explicitly credited to KDA, EDA's own contribution being a second,
independently addressed erase factor. To avoid overclaiming on
\textsc{DeltaProduct}'s behalf: its Appendix~B.4 Eq.~6 does place a per-channel
$\mathrm{diag}(w)$ inside a product at arbitrary $n_h$, but the factors are
asymmetric RWKV-7 matrices rather than Householders, the construction is offered
explicitly as buying expressivity ``at the cost of stability'', and it is used
for a theory result with no experiments; \textsc{DeltaProduct}'s own gated
variant uses a \emph{scalar} gate outside the product. We contribute no
operator, no kernel and no theorem --- only a controlled measurement.

%%%%%%%%%%%%%%%%%%%%%%%%%%%%%%%%%%%%%%%%%%%%%%%%%%%%%%%%%%%%%%%%%%%%%%%%%%%%%%%
\section{Conclusion}
\label{sec:conclusion}

In a gated delta-rule node, the write-strength range $\beta$ and the Householder
arity $R$ are not interchangeable routes to expressivity. Across 48 paired cells
on two non-solvable group word problems, widening $\beta$ produced large gains
in every position band where the compared arms were above chance --- $25.7\times$
to $124.7\times$ the gain from doubling arity --- at zero parameter cost, against
arity's $+40.3\%$. The $\beta \times R$ interaction was positive in five of eight
de-nested bands. The determinant argument of Section~\ref{sec:det}, which is not
ours, establishes independently of any experiment that a strict-$\beta$ node such
as Kimi~K3's KDA layer cannot obtain a negative-determinant transition at any
arity.

We are not able to claim more than that, and we would rather say so than imply
otherwise. Four things stand between this measurement and a result. Our
pre-registered mechanism is falsified and the fallback we offered is falsified
too, so we have no mechanism. Our arity arms are not parameter-matched, and the
bound we had hoped to place on that confound turns out to be uninformative
because both strict arms are at chance. Only one of eight band-level interactions
survives multiple-comparison correction at $n=6$, and none of the ten cells on
the prefix-averaged metric does. And an asymmetric optimization failure means
what we measured may be trainability rather than expressivity. Each has a
specified, cheap fix --- the \texttt{R1-P} capacity control, more bundles, an
$R=3$ arm, and per-arm learning-rate tuning with checkpoint retention. What we
can defend is narrower than where we started: one lever is free and appears to
matter a great deal, the other is expensive and appears to matter little, and we
do not yet know why or whether it survives at scale.

%%%%%%%%%%%%%%%%%%%%%%%%%%%%%%%%%%%%%%%%%%%%%%%%%%%%%%%%%%%%%%%%%%%%%%%%%%%%%%%
% Manual bibliography -- no bibtex pass.
%%%%%%%%%%%%%%%%%%%%%%%%%%%%%%%%%%%%%%%%%%%%%%%%%%%%%%%%%%%%%%%%%%%%%%%%%%%%%%%
\small
\begin{thebibliography}{99}

\bibitem{dao2024mamba2}
T.~Dao and A.~Gu.
\newblock Transformers are SSMs: Generalized models and efficient algorithms
through structured state space duality.
\newblock \emph{ICML}, 2024. arXiv:2405.21060.

\bibitem{li2026eda}
X.~Li, C.~Zhang, H.~Luo, X.~Lin, Z.~Wang, Z.~Qiu, Y.~Mao, L.~Chen, M.~Yuan,
M.~Sun, H.~Jiang, S.~Zhang, R.~Men, W.~Hu, G.~Cheng, B.~Zheng, D.~Liu, and
J.~Zhou.
\newblock Erase-then-Delta attention: Decoupling erase and write addresses in
delta-rule linear attention.
\newblock arXiv:2606.26560, 2026.

\bibitem{grazzi2024negative}
R.~Grazzi, J.~Siems, A.~Zela, J.~K.~H.~Franke, F.~Hutter, and M.~Pontil.
\newblock Unlocking state-tracking in linear RNNs through negative eigenvalues.
\newblock \emph{ICLR}, 2025. arXiv:2411.12537.

\bibitem{kimi2025linear}
Kimi Team.
\newblock Kimi Linear: An expressive, efficient attention architecture.
\newblock arXiv:2510.26692, 2025.

\bibitem{kimi2026k3}
Kimi Team.
\newblock Kimi K3: Open frontier intelligence.
\newblock arXiv:2607.24653, 2026.

\bibitem{schlag2021fwp}
I.~Schlag, K.~Irie, and J.~Schmidhuber.
\newblock Linear transformers are secretly fast weight programmers.
\newblock \emph{ICML}, 2021. arXiv:2102.11174.

\bibitem{siems2025deltaproduct}
J.~Siems, T.~Carstensen, A.~Zela, F.~Hutter, M.~Pontil, and R.~Grazzi.
\newblock DeltaProduct: Improving state-tracking in linear RNNs via Householder
products.
\newblock \emph{NeurIPS}, 2025. arXiv:2502.10297.

\bibitem{yang2024deltanet}
S.~Yang, B.~Wang, Y.~Zhang, Y.~Shen, and Y.~Kim.
\newblock Parallelizing linear transformers with the delta rule over sequence
length.
\newblock \emph{NeurIPS}, 2024. arXiv:2406.06484.

\bibitem{yang2024gated}
S.~Yang, J.~Kautz, and A.~Hatamizadeh.
\newblock Gated delta networks: Improving Mamba2 with delta rule.
\newblock \emph{ICLR}, 2025. arXiv:2412.06464.

\end{thebibliography}

\end{document}
